%-----------------------------------------------------------------------------------------------------------------------------------------------%
%   The MIT License (MIT)
%
%   Copyright (c) 2021 Jitin Nair
%
%   Permission is hereby granted, free of charge, to any person obtaining a copy
%   of this software and associated documentation files (the "Software"), to deal
%   in the Software without restriction, including without limitation the rights
%   to use, copy, modify, merge, publish, distribute, sublicense, and/or sell
%   copies of the Software, and to permit persons to whom the Software is
%   furnished to do so, subject to the following conditions:
%
%   THE SOFTWARE IS PROVIDED "AS IS", WITHOUT WARRANTY OF ANY KIND, EXPRESS OR
%   IMPLIED, INCLUDING BUT NOT LIMITED TO THE WARRANTIES OF MERCHANTABILITY,
%   FITNESS FOR A PARTICULAR PURPOSE AND NONINFRINGEMENT. IN NO EVENT SHALL THE
%   AUTHORS OR COPYRIGHT HOLDERS BE LIABLE FOR ANY CLAIM, DAMAGES OR OTHER
%   LIABILITY, WHETHER IN AN ACTION OF CONTRACT, TORT OR OTHERWISE, ARISING FROM,
%   OUT OF OR IN CONNECTION WITH THE SOFTWARE OR THE USE OR OTHER DEALINGS IN
%   THE SOFTWARE.
%
%-----------------------------------------------------------------------------------------------------------------------------------------------%

%----------------------------------------------------------------------------------------
%   DOCUMENT DEFINITION
%----------------------------------------------------------------------------------------

% article class because we want to fully customize the page and not use a cv template
\documentclass[a4paper,10pt]{article}

%----------------------------------------------------------------------------------------
%   FONT
%----------------------------------------------------------------------------------------

% % fontspec allows you to use TTF/OTF fonts directly
% \usepackage{fontspec}
% \defaultfontfeatures{Ligatures=TeX}

% % modified for ShareLaTeX use
% \setmainfont[
% SmallCapsFont = Fontin-SmallCaps.otf,
% BoldFont = Fontin-Bold.otf,
% ItalicFont = Fontin-Italic.otf
% ]
% {Fontin.otf}

%----------------------------------------------------------------------------------------
%   PACKAGES
%----------------------------------------------------------------------------------------
\usepackage{url}
\usepackage{parskip}

%other packages for formatting
\RequirePackage{color}
\RequirePackage{graphicx}
\usepackage[usenames,dvipsnames]{xcolor}
\usepackage[margin=0.36in]{geometry}

%tabularx environment
\usepackage{tabularx}

%for lists within experience section
\usepackage{enumitem}

% centered version of 'X' col. type
\newcolumntype{C}{>{\centering\arraybackslash}X}

%to prevent spillover of tabular into next pages
\usepackage{supertabular}
\usepackage{tabularx}
\newlength{\fullcollw}
\setlength{\fullcollw}{0.47\textwidth}

%custom \section
\usepackage{titlesec}
\usepackage{multicol}
\usepackage{multirow}

%CV Sections inspired by:
%http://stefano.italians.nl/archives/26
\titleformat{\section}{\Large\scshape\raggedright}{}{0em}{}[\titlerule]
\titlespacing{\section}{0pt}{8pt}{4pt}

\setlist[itemize]{
    nosep,
    leftmargin=1.2em,
    itemsep=2pt,
    topsep=2pt,
    parsep=0pt,
    partopsep=0pt
}

\newcommand{\jobheader}[2]{%
    \par\noindent\textbf{#1}\hfill #2\par
}
\newcommand{\jobrole}[1]{%
    \textit{#1}\par
}
\newcommand{\jobgap}{\par\addvspace{8pt}}

%for publications
% \usepackage[style=authoryear,sorting=ynt, maxbibnames=2]{biblatex}

%Setup hyperref package, and colours for links
\usepackage[unicode, draft=false]{hyperref}
\definecolor{linkcolour}{rgb}{0,0.2,0.6}
\hypersetup{colorlinks,breaklinks,urlcolor=linkcolour,linkcolor=linkcolour}
% \addbibresource{citations.bib}
% \setlength\bibitemsep{1em}

%for social icons
\usepackage{fontawesome5}

%debug page outer frames
%\usepackage{showframe}

%----------------------------------------------------------------------------------------
%   BEGIN DOCUMENT
%----------------------------------------------------------------------------------------
\begin{document}
\setlength{\parskip}{0pt}
\setlength{\parindent}{0pt}
\linespread{0.97}\selectfont
\emergencystretch=2em

% non-numbered pages
\pagestyle{empty}

%----------------------------------------------------------------------------------------
%   TITLE
%----------------------------------------------------------------------------------------

{\centering
{\Huge HARSH PRATAP SINGH}\\[4pt]
\href{https://github.com/harsh-ps-2003}{\raisebox{-0.05\height}\faGithub\ harsh-ps-2003} \ $|$ \
\href{https://harsh-ps-2003.github.io/}{\raisebox{-0.05\height}\faGlobe \ personal blogs} \ $|$ \
\href{mailto:harshpratapsingh8210@gmail.com}{\raisebox{-0.05\height}\faEnvelope \ harshpratapsingh8210@gmail.com}\par
}

%----------------------------------------------------------------------------------------
% EXPERIENCE SECTIONS
%----------------------------------------------------------------------------------------

\section{Professional Experience}

\jobheader{Sqwish AI{\normalfont, Cambridge, England}}{September 2026 -- Present}
\jobrole{Founding Member of Technical Staff}
Building low-latency contextual bandit inference for \href{https://routeworks.github.io/}{Sqwish Router}, enabling sub-50ms routing decisions for AI agents.

\jobgap
\jobheader{Cloud Native Computing Foundation (CNCF)}{March 2025 -- May 2025}
\jobrole{Linux Foundation Mentee $|$ Member of Istio's Networking Working Group}
\begin{itemize}
    \item[--] \href{https://docs.google.com/document/d/19vMr22inhLAPjWcVJkZWiG_GY6ZG3PLxG9QqAmIlnIc/edit?tab=t.0\#heading=h.5d4q82tdn6qhq}{Designed} and \href{https://github.com/istio/ztunnel/pull/1507}{implemented} \textbf{mTLS}-secured metrics access in Istio Ambient's Rust-based ztunnel, introducing a self-proxy inbound listener that leveraged workload identities and \textbf{HBONE} tunneling to secure observability traffic.
    \item[--] Integrated the feature with ztunnel's existing certificate management and inbound proxy infrastructure, and added \mbox{end-to-end} tests validating authenticated metrics scraping and mTLS telemetry.
\end{itemize}

\jobgap
% \jobheader{Ethereum Foundation}{June 2025 -- November 2025}
% \jobrole{Protocol Engineering Grantee -- EF Protocol Consensus Team}
% Implemented \href{https://arxiv.org/pdf/2405.00549}{Fast Confirmation Rule} in Ethereum Clients, providing block confirmation in \textbf{12-24 seconds vs traditional 13-19 minutes} on the Ethereum Blockchain.
%
% \jobgap
\jobheader{uv}{}
\jobrole{Open Source Contributor}
\begin{itemize}
    \item[--] Added GitLab trusted publisher support with \textbf{secure OIDC token discovery} for passwordless PyPI auth - \href{https://github.com/astral-sh/uv/pull/15583}{PR}
    \item[--] Improved packaging UX with compressed package size reporting - \href{https://github.com/astral-sh/uv/pull/15531}{PR}, authoritative local-over-global pin resolution - \href{https://github.com/astral-sh/uv/pull/15473}{PR}, and a context-aware warning for conflicting VIRTUALENV setups - \href{https://github.com/astral-sh/uv/pull/15467}{PR}
\end{itemize}

\jobgap
\jobheader{Jenkins}{May 2023 -- Present}
\jobrole{Open Source Maintainer $|$ Google Summer of Code}
\begin{itemize}
    \item[--] Nurtured Google Summer of Code contributors and worked with them to make Jenkins more AI native
    \item[--] Migrated the \href{https://www.jenkins.io/blog/2023/08/24/gitlab-plugin-modernization-report/}{Jenkins GitLab Plugin} from RESTEasy to GitLab4J-API with reverse-proxy support and integration tests as a \textbf{Google Summer of Code 2023 mentee}
\end{itemize}

%----------------------------------------------------------------------------------------
%   RESEARCH PUBLICATIONS
%----------------------------------------------------------------------------------------

\section{Research Publications}
\jobheader{Matching the wavefield is not matching the CPML solver}{\textit{\href{https://ai4sciencecommunity.github.io/neurips26.html}{NeurIPS 2026 AI for Science Workshop}}}
\textit{\href{https://home.iitk.ac.in/~dghosal/Home.html}{Prof. Dibakar Ghosal} and \href{https://www.cse.iitk.ac.in/users/subhajit/}{Subhajit Roy}} $|$ \href{https://harsh-ps-2003.github.io/writes/2026-01-02-the-chronicles-of-training-fbpinns/}{blog} $|$ \href{https://harsh-ps-2003.github.io/about/neurIPS_ai4science.pdf}{paper}
\begin{itemize}
    \item[--] Improved full-waveform inversion by showing that pressure-only \textbf{FNO}, \textbf{DeepONet}, and \textbf{SIREN} models can match the receivers and still miss the \textbf{CPML} solver. The \textbf{FNO} takes \textbf{11 harmful inversion steps out of 100} at \textbf{0.10\%} receiver error, so a close match at the receivers does not mean the subsurface model got better.
    \item[--] Used multiscale \href{https://arxiv.org/pdf/2107.07871}{\textbf{FBPINNs}} for successful forward modeling and inversion in complex \textbf{2D acoustic} geological settings.
    \item[--] Executed forward and inverse inversion end to end on dual V100s, reaching \textbf{15x} faster runs and \textbf{56\% less memory per GPU} at 75 subdomains, so it can grow past \textbf{1000 subdomains}.
\end{itemize}

\jobgap
\jobheader{Multi-GPU Hash Tables under UVM Oversubscription}{\textit{Preliminary, \href{https://www.cse.iitk.ac.in/users/swarnendu/files/papers/gpu-ovs-ds-oopsla26.pdf}{OOPSLA 2026}}}
\textit{\href{https://www.cse.iitk.ac.in/users/swarnendu/}{Prof. Swarnendu Biswas}}
\begin{itemize}
    \item[--] Optimized multi-GPU hash tables so the working set no longer has to fit in GPU memory, and at \textbf{9.1x} oversubscription they still match in-memory speed, \textbf{365M ops/s} per GPU against a \textbf{355M} ceiling.
    \item[--] Developed a \textbf{3-phase radix-routing scheduler} so each key stays on its owning GPU and pages stop thrashing across devices.
    \item[--] Achieved \textbf{76x} less inter-GPU traffic (\textbf{24.4 GB to 320 MB}) and \textbf{7.3x} throughput (\textbf{24.68M to 180.86M ops/s}), so the time goes to the lookup instead of the bus.
\end{itemize}

% \textbf{Undergraduate Project :} \hfill  \\
% Advisor: \textit{Prof. Yatinder Nath Singh (Department of Electrical Engineering)} \\
% % {\textit{Skills - Rust, gRPC, libp2p, Consistent Hashing}} \\[1pt]
% An asynchronous multidimensional distributed hash table implementation based on \href{https://pdos.csail.mit.edu/papers/chord:sigcomm01/chord_sigcomm.pdf}{Chord} DHT - {\href{https://github.com/harsh-ps-2003/pikachu}{Pikachu}}
% \begin{itemize}[
%     nosep, % No extra vertical space between items
%     after=\strut, % Ensures consistent baseline for all items
%     leftmargin=2em, % Adjust this value to control the indentation of the entire list
%     itemsep=3pt, % Space between individual items
%     label={--} % This sets the bullet symbol to a hyphen
% ]
%     \item[--] Engineered an asynchronous, in-memory \textbf{DHT} based on \href{https://pdos.csail.mit.edu/papers/chord:sigcomm01/chord_sigcomm.pdf}{Chord} with multi-dimensional keyspace support, \textbf{finger tables} and recursive lookup logic, achieving \textbf{average O(log N) hop count for key resolution}
%     \item[--] Built a resilient \textbf{routing maintenance} subsystem with \textbf{periodic stabilization}, successor/predecessor checks, and finger-table repairs to ensure network consistency under churn
%     \item[--] Engineered \textbf{connection-pooled gRPC layer with backpressure handling and bounded buffers}, reducing connection churn and ensuring stable performance during \textbf{high-volume key transfers }
%     \item[--] Introduced \textbf{streaming handoff APIs} for efficient ownership transfer of large key ranges during join/leave, \textbf{reducing tail latencies} and avoiding head-of-line blocking
%     \item[--] Implemented robust server \textbf{lifecycle management} with \textbf{readiness signalling}, \textbf{bounded startup timeouts}, and \textbf{graceful shutdown} for clean drain of in-flight requests
% \end{itemize}

% \textbf{Undergraduate Project :} \hfill  \\
% Advisor: \textit{Prof. Hamim Zafar (Department of Biological Sciences and Bioengineering)} \\
% Mosaic single-cell data integration using deep generative models
% % {\textit{Skills - Python, PyTorch, Transformers, U-Net}} 
% \begin{itemize}[
%     nosep, % No extra vertical space between items
%     after=\strut, % Ensures consistent baseline for all items
%     leftmargin=2em, % Adjust this value to control the indentation of the entire list
%     itemsep=3pt, % Space between individual items
%     label={--} % This sets the bullet symbol to a hyphen
% ]
%     \item[--] Explored various deep generative models like scMODAL, MaxFuse, Harmony, and developed an improved scMGAT to integrate mosaic single-cell multi-omics data (RNA+ADT), replacing unstable GAN baselines.
%     \item[--] Built a \textbf{dual-branch autoencoder} with a \textbf{geometry-preserving loss} (cosine-kernel similarity + progressive scheduling) to retain local and global manifold structure in latent embeddings, improving downstream graph quality
%     \item[--] Introduced \textbf{early-stage supervised contrastive learning} within a multi-head GAT to explicitly align same-cell representations across modalities, significantly improving biological separability and cross-modal label transfer
%     \item Achieved state-of-the-art integration performance on CITE-seq PBMC data, \textbf{improving clustering and cross-modal alignment with ARI = 0.889}, \textbf{NMI = 0.886}, and RNA$\leftrightarrow$ADT label transfer accuracy from \textbf{$\sim$60\% to $\sim$82\%}, alongside substantially reduced FOSCTTM alignment error
% \end{itemize}

% \textbf{Undergraduate Project :} \hfill  \\
% Advisor: \textit{Prof. Arnab Bhattacharya (Department of Computer Science)} \\
% A generative model to produce clinically accurate synthetic images across various medical modalities - \href{https://github.com/harsh-ps-2003/CS787}{Text2Med}
% % {\textit{Skills - Python, PyTorch, Transformers, U-Net}} 
% \begin{itemize}[
%     nosep, % No extra vertical space between items
%     after=\strut, % Ensures consistent baseline for all items
%     leftmargin=2em, % Adjust this value to control the indentation of the entire list
%     itemsep=3pt, % Space between individual items
%     label={--} % This sets the bullet symbol to a hyphen
% ]
%     \item[--] Developed a \textbf{multi-modal text-to-image generative model} for producing synthetic medical images across modalities using \textbf{diffusion}-based architectures with \textbf{cross-attention mechanisms}
%     \item[--] Engineered end-to-end training pipeline supporting image-text datasets with \textbf{CSV-based data ingestion}, \textbf{medically-trained BERT} tokenization for multi-modal inputs, and \textbf{U-Net architecture} with \textbf{progressive noise scheduling} for diffusion model training
%     \item[--] Built \textbf{two-stage RLHF pipeline} to iteratively refine synthetic medical images utilizing a \textbf{custom Vision Transformer} model built with \textbf{multi-head self-attention} and \textbf{dropout regularization}, leveraging radiologist ratings to serve as an automated quality-assessment reward model
%     \item [--] Optimized \textbf{Stable Diffusion} fine-tuning with \textbf{mixed precision}, \textbf{gradient checkpointing}, \textbf{cosine annealing learning rate scheduling}, \textbf{EMA model averaging}, and \textbf{SNR-weighted loss computation} with gamma parameter for noise prediction training
%     \item [--] Achieved significant memory reduction via \textbf{8-bit Adam optimization}, \textbf{attention slicing}, \textbf{VAE tiling}, and \textbf{gradient accumulation} 
%     \item [--] Developed evaluation metrics including \textbf{FID}, \textbf{Inception Score}, \textbf{SSIM}, and custom medical quality metrics
%     \item[--] Implemented \textbf{multi-GPU distributed training with Accelerate} framework, automatic device detection and allocation, \textbf{gradient synchronization across processes}, \textbf{checkpoint management} with automatic resumption, \textbf{mixed precision} coordination, and \textbf{batch size scaling}
% \end{itemize}

% \textbf{Undergraduate Project :} \hfill  \\
% Advisor: \textit{Prof. Adithya Vadapalli (Department of Computer Science)} \\
% % {\textit{Skills - Rust, Cryptography}} \\[1pt]
% A privacy-preserving discovery service \href{https://signal.org/blog/building-faster-oram/}{inspired by Signal} - \href{https://github.com/harsh-ps-2003/rumi}{Rumi}
% \begin{itemize}[
%     nosep, % No extra vertical space between items
%     after=\strut, % Ensures consistent baseline for all items
%     leftmargin=2em, % Adjust this value to control the indentation of the entire list
%     itemsep=3pt, % Space between individual items
%     label={--} % This sets the bullet symbol to a hyphen
% ]
%     \item[--] Engineered a privacy-preserving contact discovery service using \textbf{double-sided blinding over elliptic curves and compressed point encodings}, a \textbf{Path ORAM}-backed store, and \textbf{zero-knowledge set-membership proofs} to ensure the server learns nothing about queries, matches, or access patterns via fixed-size responses and uniform access paths and explicit secure wiping of ORAM blocks, stash, and internal buffers
%     \item[--] Designed prefix-based bucket selection paired with fixed dummy accesses and constant response shaping on top of Path ORAM for efficient server-side filtering while maintaining indistinguishability; tunable privacy/performance balance
%     \item[--] Benchmarked ORAM performance across read/write, sequential/random patterns, and \textbf{variable payload sizes}
% \end{itemize}

%----------------------------------------------------------------------------------------
%   EDUCATION
%----------------------------------------------------------------------------------------

\section{Education}

\noindent
Bachelor's in Earth Science (major), Computer Science and Philosophy (minors) -- \textit{Indian Institute of Technology, Kanpur}
\end{document}